\documentclass[10pt, aspectratio=169]{beamer}
% Preámbulo modular para presentaciones Beamer (Modelación Estadística)
% Diseño y paleta institucional del Tecnológico de Monterrey

\documentclass[aspectratio=169,xcolor={dvipsnames,table}]{beamer}

\usepackage[utf8]{inputenc}
\usepackage[spanish,mexico]{babel}
\usepackage{amsmath,amssymb,amsfonts,mathtools}
\usepackage{graphicx}
\usepackage{listings}
\usepackage{booktabs}
\usepackage{tikz}
\usetikzlibrary{matrix,arrows,decorations.pathmorphing,shapes.geometric,positioning}

% Paleta Institucional Tecnológico de Monterrey
\definecolor{TecAzul}{HTML}{0039A6}       % Azul TEC: color primario institucional
\definecolor{TecAzulOscuro}{HTML}{1A2E51} % Azul marino oscuro
\definecolor{TecRojo}{HTML}{EC2661}       % Rojo de acento (paleta miTec)
\definecolor{TecGris}{HTML}{646464}       % Gris neutro para comentarios y subtítulos
\definecolor{TecFondoBloque}{HTML}{F4F6F9}% Fondo suave para bloques

% Configuración de Tema Beamer
\usetheme{Boadilla}
\usecolortheme[named=TecAzulOscuro]{structure}
\setbeamercolor{alerted text}{fg=TecRojo}
\setbeamercolor{palette primary}{bg=TecAzulOscuro,fg=white}
\setbeamercolor{palette secondary}{bg=TecAzul,fg=white}
\setbeamercolor{palette tertiary}{bg=TecAzulOscuro!80!black,fg=white}
\setbeamercolor{title}{fg=TecAzulOscuro,bg=TecFondoBloque}
\setbeamercolor{frametitle}{fg=TecAzulOscuro,bg=TecFondoBloque}

% Bloques personalizados elegantes
\setbeamercolor{block title}{bg=TecAzulOscuro,fg=white}
\setbeamercolor{block body}{bg=TecFondoBloque,fg=black}
\setbeamercolor{block title alerted}{bg=TecRojo,fg=white}
\setbeamercolor{block body alerted}{bg=TecRojo!8,fg=black}
\setbeamercolor{block title example}{bg=TecAzul,fg=white}
\setbeamercolor{block body example}{bg=TecAzul!8,fg=black}

% Redondear bordes y añadir sombra sutil
\setbeamertemplate{blocks}[rounded][shadow=true]
\setbeamertemplate{navigation symbols}{}

% Pie de página limpio con número de diapositiva
\setbeamertemplate{footline}{
  \leavevmode%
  \hbox{%
  \begin{beamercolorbox}[wd=.7\paperwidth,ht=2.25ex,dp=1ex,left]{author in head/foot}%
    \hspace*{2ex}\insertshortauthor~~(\insertshortinstitute) \hfill \insertshorttitle
  \end{beamercolorbox}%
  \begin{beamercolorbox}[wd=.3\paperwidth,ht=2.25ex,dp=1ex,right]{date in head/foot}%
    \insertshortdate{}\hspace*{2em}
    \insertframenumber{} / \inserttotalframenumber\hspace*{2ex}
  \end{beamercolorbox}}%
  \vskip0pt%
}

% Configuración de Listings de Python para visualización pedagógica de código
\lstset{
    language=Python,
    basicstyle=\ttfamily\footnotesize,
    keywordstyle=\color{TecAzulOscuro}\bfseries,
    stringstyle=\color{TecRojo},
    commentstyle=\color{TecGris}\itshape,
    showstringspaces=false,
    breaklines=true,
    frame=lines,
    rulecolor=\color{TecAzulOscuro!30},
    backgroundcolor=\color{TecFondoBloque},
    aboveskip=0.5em,
    belowskip=0.5em,
    literate={á}{{\'a}}1 {é}{{\'e}}1 {í}{{\'i}}1 {ó}{{\'o}}1 {ú}{{\'u}}1
             {Á}{{\'A}}1 {É}{{\'E}}1 {Í}{{\'I}}1 {Ó}{{\'O}}1 {Ú}{{\'U}}1
             {ñ}{{\~n}}1 {Ñ}{{\~N}}1 {¿}{{?`}}1 {¡}{{!`}}1
}

% Metadatos institucionales por defecto
\title[Modelación Estadística]{Modelación Estadística para Ciencia de Datos e Ingeniería}
\author[J. Castillo Colmenares]{Juliho David Castillo Colmenares}
\institute[Tec de Monterrey]{Tecnológico de Monterrey \\ Escuela de Ingeniería y Ciencias}
\date{\today}

% Comandos y macros estadísticas compartidas para las presentaciones Beamer (Metropolis)

% Conjuntos numéricos básicos
\newcommand{\R}{\mathbb{R}}
\newcommand{\N}{\mathbb{N}}
\newcommand{\Q}{\mathbb{Q}}
\newcommand{\Z}{\mathbb{Z}}
\newcommand{\set}[1]{\left\{ #1 \right\}}
\newcommand{\abs}[1]{\left|#1\right|}
\newcommand{\norm}[1]{\left\| #1 \right\|}

% Comandos estadísticos y probabilísticos del curso
\newcommand{\Var}{\operatorname{Var}}
\newcommand{\cov}{\operatorname{Cov}}
\newcommand{\corr}{\rho}
\newcommand{\comb}[2]{\begin{pmatrix} #1 \\ #2 \end{pmatrix}}
\newcommand{\s}{\sigma}
\newcommand{\particion}{\mathfrak{P}}
\newcommand{\card}[1]{\#\left( #1 \right)}
\newcommand{\seq}[3]{\set{#1}_{#2}^{#3}}
\newcommand{\E}{\mathbb{E}}
\newcommand{\Prob}{\mathbb{P}}

% Entornos pedagógicos y cajas en español académico natural
\newenvironment{discusion}{%
  \begin{alertblock}{Pregunta de Discusión}%
}{%
  \end{alertblock}%
}

\newenvironment{reflexion}{%
  \begin{alertblock}{Reflexión para el Aula}%
}{%
  \end{alertblock}%
}

\newenvironment{paradoja}{%
  \begin{alertblock}{Paradoja de Exploración}%
}{%
  \end{alertblock}%
}

\newenvironment{motivacion}{%
  \begin{exampleblock}{Motivación: Ciencia de Datos en la Práctica}%
}{%
  \end{exampleblock}%
}

\newenvironment{retoclase}{%
  \begin{block}{Reto Rápido en Equipo}%
}{%
  \end{block}%
}


\title[Regresión sobre Datos Simulados]{Sección 09.04: Regresión Lineal sobre Datos Simulados}
\subtitle{De los Parámetros Poblacionales Verdaderos al Ajuste Muestral}
\author[J. Castillo Colmenares]{Juliho Castillo Colmenares}
\institute{Tecnológico de Monterrey}
\date{\vspace{-1.2cm}}

\begin{document}

% Slide 1: Portada
\begin{frame}[plain]
  \titlepage
\end{frame}

% Slide 2: Hoja de Ruta
\begin{frame}{Hoja de Ruta --- Capítulo 09: Regresiones Lineales y Múltiples}
  \scriptsize
  ¿Dónde nos encontramos en el desarrollo modular del capítulo?
  \begin{itemize}\itemsep=0.02em
    \item \textbf{09.01} Correlación como Premisa de la Regresión
    \item \textbf{09.02} Introducción a la Regresión Lineal
    \item \textbf{09.03} Matemáticas de la Regresión: Derivación MCO y $R^2$
    \item \textbf{09.04} Regresión Lineal sobre Datos Simulados \emph{(Hoy)}
    \item \textbf{09.05} Coeficientes Óptimos, Pruebas $t$/$F$ y RSE
    \item \textbf{09.06} Implementación en Python con \texttt{statsmodels}
    \item \textbf{09.07} Regresión Múltiple, Ecuación Normal y Ridge/Lasso
    \item \textbf{09.08} Validación de Modelos y $k$-fold Cross-Validation
    \item \textbf{09.09} Diagnóstico de Residuos y Multicolinealidad (VIF)
    \item \textbf{09.10} Regresión con \texttt{scikit-learn}
    \item \textbf{09.11} Variables Categóricas y Variables Muda
    \item \textbf{09.12} Transformaciones No Lineales y Regresión Polinomial
  \end{itemize}
  \pause
  \vspace{-0.05cm}
  \begin{block}{Objetivo de la Sesión}
    Aplicar las fórmulas derivadas en la Sección 09.03 a un conjunto de datos simulado con parámetros poblacionales conocidos, comparar el ajuste MCO contra la verdad poblacional, y descubrir por qué la descomposición exacta $\text{SCT}=\text{SCR}+\text{SCD}$ requiere específicamente la recta ajustada, no cualquier recta.
  \end{block}
\end{frame}

% Slide 3: Motivación
\begin{frame}{¿Por Qué Simular Datos con Parámetros Conocidos?}
  \footnotesize
  \begin{motivacion}
    En la práctica nunca conocemos los parámetros poblacionales verdaderos $\alpha,\beta$: solo observamos datos y estimamos $\hat\alpha,\hat\beta$. Al \textbf{simular} datos, invertimos el problema: fijamos $\alpha,\beta,\sigma$ verdaderos, generamos observaciones ruidosas, y verificamos qué tan bien MCO recupera la verdad que nosotros mismos construimos.
  \end{motivacion}

  \vspace{0.15cm}
  \pause
  Escribimos el modelo poblacional con su componente de error: $Y=\alpha+\beta X+K$, donde $K\sim N(0,\sigma^2)$.
\end{frame}

% Slide 4: Diseño del experimento
\begin{frame}{Diseño del Experimento de Simulación}
  \footnotesize
  \begin{block}{Parámetros Poblacionales Verdaderos}
    $$\alpha=2.0, \qquad \beta=0.3, \qquad \sigma=0.5$$
  \end{block}
  \pause

  \vspace{0.1cm}
  \begin{alertblock}{Generación de Datos}
    $X\sim N(1.5,\,2.5^2)$, con $n=100$ observaciones. Para cada $X_i$, generamos $Y_i=\alpha+\beta X_i+K_i$ con $K_i\sim N(0,0.5^2)$: el valor "actual" observado, distinto del valor "predicho" por la recta poblacional pura.
  \end{alertblock}
\end{frame}

% Slide 5: SST, SCR, SCD
\begin{frame}{Suma de Cuadrados: Total, de Regresión, y de Diferencia}
  \footnotesize
  \begin{block}{Las Tres Cantidades}
    $$\text{SCT}=\sum(Y_i-\bar Y)^2, \quad \text{SCR}=\sum(Y_{e,i}-\bar Y)^2, \quad \text{SCD}=\sum(Y_i-Y_{e,i})^2$$
  \end{block}
  \pause

  \vspace{0.1cm}
  \begin{alertblock}{Interpretación}
    SCR es la variabilidad \emph{explicada} por el modelo; SCD es la variabilidad \emph{no explicada}. Entre mayor la proporción SCR:SCT, mejor el modelo --- esto es exactamente $R^2$.
  \end{alertblock}
\end{frame}

% Slide 6: La sutileza de esta sección
\begin{frame}{Una Sutileza Importante: ¿Qué Recta Usar?}
  \footnotesize
  \begin{alertblock}{Pregunta}
    La Sección 09.03 demostró $\text{SCT}=\text{SCR}+\text{SCD}$ para la recta \textbf{ajustada por MCO}. ¿Se cumple esta identidad si en su lugar usamos la recta poblacional \textbf{verdadera} ($\alpha=2.0,\beta=0.3$) para calcular $Y_{e,i}$?
  \end{alertblock}
  \pause

  \vspace{0.1cm}
  \begin{block}{Respuesta Adelantada}
    No exactamente --- la demostración de la Sección 09.03 depende de que $\sum e_i=0$ y $\sum X_ie_i=0$, propiedades garantizadas por las ecuaciones normales \textbf{solo} para la recta que minimiza $S(\alpha,\beta)$ en esta muestra específica.
  \end{block}
\end{frame}

% Slide 7: Lab Python (Bloque 1)
\begin{frame}[fragile]{Laboratorio en Python: Simulación del Modelo Poblacional}
  \scriptsize
  Generación de datos con parámetros verdaderos conocidos (\texttt{09.04\_regression\_on\_simulated\_data.py}):
  {\lstset{basicstyle=\fontsize{3.6pt}{4.3pt}\selectfont\ttfamily,aboveskip=0.05em,belowskip=0.05em}
  \lstinputlisting[firstline=17, lastline=27]{../../code/09_regresiones/09.04_regression_on_simulated_data.py}}
\end{frame}

% Slide 8: Lab Python (Bloque 2)
\begin{frame}[fragile]{Laboratorio en Python: Ajuste MCO vs. Verdad Poblacional}
  \scriptsize
  Comparación de $\hat\alpha,\hat\beta$ contra $\alpha,\beta$ verdaderos, y verificación exacta de SCT=SCR+SCD:
  {\lstset{basicstyle=\fontsize{3.2pt}{3.9pt}\selectfont\ttfamily,aboveskip=0.05em,belowskip=0.05em}
  \lstinputlisting[firstline=30, lastline=45]{../../code/09_regresiones/09.04_regression_on_simulated_data.py}}
\end{frame}

% Slide 9: Lab Python (Bloque 3)
\begin{frame}[fragile]{Laboratorio en Python: Recta Ajustada vs. Recta Oráculo}
  \scriptsize
  La descomposición exacta solo se cumple para la recta ajustada, no para la recta poblacional verdadera:
  {\lstset{basicstyle=\fontsize{3.0pt}{3.7pt}\selectfont\ttfamily,aboveskip=0.05em,belowskip=0.05em}
  \lstinputlisting[firstline=48, lastline=68]{../../code/09_regresiones/09.04_regression_on_simulated_data.py}}
\end{frame}

% Slide 10: Salida Terminal
\begin{frame}[fragile]{Laboratorio en Python: Salida en Terminal}
  \scriptsize
  Salida estándar generada al ejecutar el script del laboratorio en Python:
  \vspace{0.02cm}
  \begin{block}{Salida Estándar en Consola}
    \fontsize{3.6pt}{4.3pt}\selectfont
    \begin{verbatim}
=== Block 1: Simulation ===
True: alpha=2.0, beta=0.3, sigma=0.5; n=100

=== Block 2: OLS Fit vs. Truth ===
Fitted:  alpha_hat=2.0166, beta_hat=0.3105
True:    alpha=2.0000,     beta=0.3000
Fitted-line: SST=102.2657, SSR=75.6532, SSE=26.6126, SSR+SSE=102.2657 (exact match)

=== Block 3: Fitted vs. Oracle Line ===
Fitted line:  SSR+SSE=102.2657 (matches SST), R^2=0.7398
Oracle line:  SSR+SSE=97.5712  (does NOT match SST=102.2657), R^2=0.6918
    \end{verbatim}
  \end{block}
\end{frame}

% Slide 11: Interpretación de resultados
\begin{frame}{Interpretación: Estimación vs. Verdad}
  \footnotesize
  \begin{itemize}\itemsep=0.12cm
    \item Los estimadores $\hat\alpha\approx2.017$, $\hat\beta\approx0.311$ están muy cerca de los verdaderos $\alpha=2.0$, $\beta=0.3$: evidencia empírica del insesgo demostrado analíticamente en la Sección 09.03. \pause
    \item El $R^2\approx0.74$ de la recta ajustada es, por construcción, el \textbf{máximo} $R^2$ alcanzable con un modelo lineal en esta muestra --- ninguna otra recta, ni siquiera la poblacional verdadera, puede superarlo. \pause
    \item Esto es consecuencia directa de que MCO minimiza SCD (y por tanto maximiza SCR) \emph{sobre esta muestra específica}, no sobre la población.
  \end{itemize}
\end{frame}

% Slide 12: Cierre
\begin{frame}{Cierre de Sección: Regresión sobre Datos Simulados}
  \begin{reflexion}
    La simulación con parámetros conocidos nos permite verificar empíricamente que MCO recupera la verdad poblacional con precisión creciente, y revela una sutileza importante: la descomposición exacta SCT=SCR+SCD --- y por tanto el $R^2$ máximo posible --- pertenece exclusivamente a la recta ajustada por mínimos cuadrados sobre la muestra observada.
  \end{reflexion}

  \vspace{0.2cm}
  \pause
  \begin{block}{Lecciones Clave}
    \begin{itemize}\itemsep=0.08cm
      \item \textbf{Simulación como verificación:} generar datos con parámetros conocidos es la forma más directa de auditar un método de estimación.
      \item \textbf{MCO es óptimo \emph{muestralmente}:} minimiza el error en \emph{esta} muestra, no necesariamente iguala a la recta poblacional.
      \item \textbf{La identidad SCT=SCR+SCD es exclusiva del ajuste MCO}, no de cualquier recta razonable.
    \end{itemize}
  \end{block}
\end{frame}

% Slide 13: Perspectiva Modular
\begin{frame}{Perspectiva Modular --- El Próximo Reto}
  \scriptsize
  \begin{retoclase}
    Hemos visto que $\hat\alpha,\hat\beta$ son estimadores puntuales, pero ¿qué tan confiables son? La Sección 09.05 responde formalmente: someteremos $\hat\beta$ a una prueba de hipótesis para verificar su significancia estadística, construiremos su intervalo de confianza, y calcularemos el Error Estándar Residual (RSE) como medida absoluta de la calidad del ajuste.
  \end{retoclase}

  \vspace{0.08cm}
  \pause
  \begin{alertblock}{Preparación Recomendada}
    \begin{itemize}\itemsep=0.06cm
      \item Repasa la distribución muestral de un estimador (Capítulo 05) y su relación con la construcción de pruebas de hipótesis.
      \item Reflexiona sobre qué significaría que $\hat\beta$ no fuera estadísticamente distinto de cero.
    \end{itemize}
  \end{alertblock}
\end{frame}

\end{document}
